\documentclass[12pt]{article}
\usepackage[margin=1in]{geometry}
\usepackage{setspace}
\doublespacing

%% Bibliography
\usepackage{natbib}
\bibliographystyle{aea}

%% Tables
\usepackage{booktabs, threeparttable, siunitx, caption}
\usepackage{tabu}
\usepackage{float}
\extrarowsep=0.1pt
\renewcommand{\arraystretch}{1.15}

%% Figures
\usepackage{graphicx}
\usepackage{subcaption}

%% Math & fonts
\usepackage{amsmath, amssymb}
\usepackage{mathpazo}

%% Cross-references
\usepackage{hyperref}
\hypersetup{colorlinks=true, linkcolor=black, citecolor=black, urlcolor=black}

%% Misc
\usepackage[dvipsnames]{xcolor}
\usepackage{microtype}

\def\sym#1{\ifmmode^{#1}\else\(^{#1}\)\fi}

% ─────────────────────────────────────────────────────────────────────────────
\begin{document}

% ── Title ────────────────────────────────────────────────────────────────────
\begin{titlepage}
\begin{center}
\vspace*{1.5cm}
{\large \textbf{Does Baseline Willingness-to-Pay Predict\\[0.3em]
Life Satisfaction Gains from Electrification?\\[0.3em]
Evidence from Rural Kenya}}\\[2em]
{\normalsize Extension of \citet{lee2019jpe}}\\[1.5em]
{\normalsize \today}
\end{center}

\vspace{2em}
\begin{abstract}
\noindent
We ask whether a household's stated willingness-to-pay (WTP) for a grid
electricity connection at baseline predicts the life satisfaction gains it
realises from connection. Using data from the randomised experiment of
\citet{lee2019jpe}, we merge a contingent-valuation WTP measure from the
2014 baseline survey with subjective wellbeing outcomes from two follow-up
rounds (2016, 2017). We estimate heterogeneous intent-to-treat (ITT) effects
on a normalised life satisfaction index and a perceptions-of-security index,
interacting the treatment indicator with a binary high-WTP variable. The
differential ITT coefficient is small (${-}0.018$ standard deviations) and
statistically insignificant in all specifications. After Benjamini--Hochberg
correction, neither outcome survives multiple-testing adjustment
($q \geq 0.95$). These results are consistent with the null average
wellbeing effects reported by \citet{lee2019jpe} and suggest that
prior expectations, as proxied by stated WTP, do not meaningfully
moderate the wellbeing impact of grid connection.
\end{abstract}
\end{titlepage}

% ── 1. Background and Motivation ─────────────────────────────────────────────
\section{Background and Motivation}
\label{sec:background}

A growing literature examines the welfare effects of household
electrification in developing countries. \citet{lee2019jpe} conduct a
randomised experiment in rural Kenya in which households are randomly
offered grid connections at varying subsidy levels. Despite large take-up
responses to price subsidies, the authors find small or null average
effects on income, employment, and subjective wellbeing.

One mechanism that could generate heterogeneous wellbeing effects even
within an experiment with null average effects is an \emph{expectations
channel}: households that valued the connection more highly ex ante---as
revealed by higher stated WTP---may have had stronger prior expectations
of benefit. If electrification meets those expectations, high-WTP
households should exhibit larger life satisfaction gains than low-WTP
households; if it fails to meet them, the reverse could hold.

This note tests the expectations channel directly by interacting the
experimental treatment indicator with a baseline WTP measure. The
specification follows \citet{lee2019jep} Appendix Tables A3--A4, which
interact grid connection with a high-SES composite index. We substitute
the SES index with a binary high-WTP indicator derived from the
contingent-valuation exercise conducted at baseline.

% ── 2. Data and Specification ────────────────────────────────────────────────
\section{Data and Specification}
\label{sec:data}

\subsection*{Data}

We use three datasets from \citet{lee2019jpe}. The baseline survey
(\texttt{demand.dta}, 2014) contains the WTP elicitation and treatment
assignment. Two follow-up surveys---Round 1 (\texttt{impacts\_1.dta},
2016) and Round 2 (\texttt{impacts\_2.dta}, 2017)---contain our outcome
variables. All three share a unique household identifier
(\texttt{reppid}), enabling a 1:1 merge. We drop households already
connected at baseline (\texttt{base\_connected}$=1$) and restrict to
adult respondents (\texttt{child}$=0$), following the original papers.

\subsection*{WTP Measure}

At baseline, each household was randomly assigned a hypothetical
connection price (\texttt{WTP\_amt} $\in \{0, 10{,}000, \ldots,
75{,}000\}$ KES) and asked whether it would accept
(\texttt{WTP\_r1}$\in\{0,1\}$). We define stated WTP as:
\begin{equation}
\label{eq:wtp}
\text{WTP-stated}_i = \texttt{WTP\_amt}_i \times
\mathbf{1}[\texttt{WTP\_r1}_i = 1].
\end{equation}
Since the elicitation price is randomly assigned, stated WTP has an
exogenous component. The sample median of WTP-stated equals zero (56.4\,\%
of households declined the offer), so \texttt{wtp\_high}$_i = 1$
effectively identifies the 43.6\,\% of households that accepted at any
positive price.

\subsection*{Outcomes}

Our primary outcome is \texttt{life\_index}, a normalised life
satisfaction index available in both R1 and R2. Our secondary---and
exploratory---outcome is \texttt{r\_crime\_index}, the perceptions-of-%
security index (JPE Table 3, row D6), available in R2 only. Both
indices are standardised to mean zero and unit standard deviation in the
control group.

\subsection*{Specification}

We estimate:
\begin{equation}
\label{eq:spec}
y_{ic} = \alpha
  + b_1\,\text{Treat}_{ic}
  + b_2\,\text{HighWTP}_{i}
  + b_3\,(\text{Treat}_{ic} \times \text{HighWTP}_{i})
  + \boldsymbol{\delta}'\mathbf{X}_{ic}
  + \varepsilon_{ic},
\end{equation}
where $\text{Treat}_{ic} = \texttt{treated\_all}$ (any subsidy arm),
$\text{HighWTP}_i = \mathbf{1}[\text{WTP-stated}_i > \text{median}]$,
and $\mathbf{X}_{ic}$ includes community controls (county, market
proximity, transformer timing, baseline connection rate, population) and
household controls (gender, age, housing quality, baseline energy spending).
For the pooled D2 specification we include a Round~2 fixed effect.

\noindent
The coefficient of interest is $b_3$: a positive and significant
estimate would indicate that high-WTP households benefit more from
treatment in terms of wellbeing. Standard errors are clustered at the
community level (\texttt{siteno}). We apply the Benjamini--Hochberg
procedure to adjust for multiple testing across specifications
\citep{benjamini1995}.

% ── 3. Results ───────────────────────────────────────────────────────────────
\section{Results}
\label{sec:results}

Figure~\ref{fig:coefplot} plots ITT coefficients separately for
low-WTP (hollow circles) and high-WTP (filled circles) households,
with 95\,\% confidence intervals. For both D2 (life satisfaction) and
D6 (security index), the two groups' estimates are close together and
straddle zero, with overlapping confidence intervals. There is no visual
evidence of a positive expectations effect.

\begin{figure}[H]
  \centering
  \includegraphics[width=0.92\textwidth]{%
    ../../outputs/figures/wellbeing_wtp_heterogeneity.pdf}
  \caption{ITT coefficients on wellbeing outcomes by baseline WTP group.
    Hollow circles: low-WTP households (WTP-stated $\leq$ median).
    Filled circles: high-WTP households (WTP-stated $>$ median).
    Horizontal bars are 95\,\% confidence intervals. Vertical dashed
    line at zero. D6 is exploratory (R2 only; does not survive FDR
    correction). See Section~\ref{sec:data} for specification details.}
  \label{fig:coefplot}
\end{figure}

Table~\ref{tab:main} presents the full regression results in the format
of \citet{lee2019jpe} Table~3. For D2 (Panel~D, row~D2), the low-WTP
ITT is $0.094^{**}$ ($p = 0.042$), consistent with the positive average
effect found in the original paper. However, the high-WTP ITT is
$0.076$ ($p = 0.127$), and the differential $\hat{b}_3 = -0.018$
($\text{SE} = 0.066$, $p = 0.787$). After Benjamini--Hochberg correction
the FDR $q$-value for the differential is $0.950$, far from conventional
significance thresholds.

For D6 (Panel~D, row~D6), neither the low-WTP nor the high-WTP ITT
is significant, and the differential is $-0.012$ ($p = 0.893$,
$q = 0.950$). Panel~E shows the mean treatment effect on the
noneconomic index (average of standardised D2 and D6): the differential
is $-0.047$ ($p = 0.486$), again insignificant.

\clearpage
\newgeometry{left=0.9in, right=0.9in, top=1in, bottom=1in}
\begin{singlespace}
%% Tab: Wellbeing Heterogeneity by Baseline WTP
%% Standalone table — \input this file into a document that loads:
%%   booktabs, threeparttable, siunitx, caption, tabu, setspace

\def\sym#1{\ifmmode^{#1}\else\(^{#1}\)\fi}

\begin{table}[htbp]
\centering
\caption*{Table 1\textemdash Heterogeneous effects of grid connection on wellbeing,
  by baseline willingness-to-pay}
\vspace{-0.5em}
\begin{threeparttable}
\begin{tabu} to \linewidth {X[20,l] X[5,c] X[5,c] X[5,c] X[5,c] X[4,c]}
\toprule
 & Control & Low-WTP & High-WTP & Treat\,$\times$ & FDR \\[-0.25em]
 & & ITT & ITT & HighWTP & $q$-value \\[-0.25em]
 & (1) & (2) & (3) & (4) & (5) \\
\midrule \\[-1.2em]

%% ── Panel D ──────────────────────────────────────────────────────────────
\multicolumn{6}{l}{\emph{Panel D: Primary Noneconomic Outcomes}} \\[0.5em]

\hspace{2.5mm} D2. Normalised life satisfaction
  & 0.000 & $0.094$\sym{**} & $0.076$ & $-0.018$ & 0.950 \\
& [1.000] & $(0.046)$ & $(0.050)$ & $(0.066)$ & \\[0.8em]

\hspace{2.5mm} D6. Perceptions of security index\tnote{$\dagger$}
  & 0.000 & $-0.048$ & $-0.060$ & $-0.012$ & 0.950 \\
& [1.000] & $(0.059)$ & $(0.073)$ & $(0.089)$ & \\[0.8em]

%% ── Panel E ──────────────────────────────────────────────────────────────
\multicolumn{6}{l}{\emph{Panel E: Mean Treatment Effects on Grouped Outcomes}} \\[0.5em]

\hspace{2.5mm} E1. Noneconomic index (D2\,+\,D6)\tnote{$\ddagger$}
  & 0.000 & $0.038$ & $-0.009$ & $-0.047$ & {---} \\
& [0.700] & $(0.047)$ & $(0.052)$ & $(0.067)$ & \\

\bottomrule
\end{tabu}

\begin{tablenotes}[flushleft]
\footnotesize
\setlength{\itemsep}{2pt}
\item \textit{Notes}: Column (1) reports the control-group mean with the
  standard deviation in brackets. Columns (2) and (3) report ITT
  coefficients from the interaction regression
  $y_{ic} = \alpha + b_1\,\text{Treat}_{ic} + b_2\,\text{HighWTP}_{i}
  + b_3\,(\text{Treat}\times\text{HighWTP})_{ic} + \boldsymbol{\delta}'\mathbf{X}_{ic}
  + \varepsilon_{ic}$,
  where column (2) is $\hat{b}_1$ (low-WTP ITT) and column (3) is
  $\hat{b}_1 + \hat{b}_3$ (high-WTP ITT); standard errors computed by
  the delta method. Column (4) reports $\hat{b}_3$ with its standard
  error in parentheses; column (5) reports the Benjamini--Hochberg
  FDR-adjusted $q$-value for column (4).
  \text{Treat} $=$ \texttt{treated\_all} (any subsidy arm).
  \text{HighWTP} $=$ 1 if stated WTP above sample median (43.6\,\%
  of sample); stated WTP defined as \texttt{WTP\_amt}\,$\times$\,
  \texttt{WTP\_r1}. D2 uses pooled R1\,+\,R2 data ($N = 4{,}157$);
  D6 and E1 use R2 only ($N = 2{,}047$--$2{,}057$).
  Controls: county, market proximity, transformer timing, baseline
  connection rate, population, gender, age, housing quality, baseline
  energy spending. Standard errors clustered at the community level.
  Asterisks indicate significance (two-tailed):
  $^{*}p<0.10$;\;$^{**}p<0.05$;\;$^{***}p<0.01$.
\item[$\dagger$] D6 is exploratory. It does not survive
  Benjamini--Hochberg correction ($q = 0.950$) and is available in
  R2 only (\texttt{r\_crime\_index}).
\item[$\ddagger$] E1 is the mean of the standardised D2 and D6
  z-scores. FDR $q$-value not reported for Panel E following
  Lee, Miguel \& Wolfram (2019).
\end{tablenotes}
\end{threeparttable}
\end{table}

\end{singlespace}
\restoregeometry
\addtocounter{table}{0}%
\renewcommand{\thetable}{\arabic{table}}%

Across all three rounds and both with and without controls, the point
estimates of $b_3$ are small and consistently negative, ranging from
$-0.071$ (R2, D2, with controls) to $+0.032$ (R1, D2, with controls).
The robustness specification using a continuous WTP measure
(WTP-stated in USD) similarly yields an insignificant interaction
($\hat{b}_3 = 0.0001$, $p = 0.647$).

% ── 4. Caveats ───────────────────────────────────────────────────────────────
\section{Caveats}
\label{sec:caveats}

\paragraph{WTP measure is binary and partially random.}
Because each household sees a single randomly assigned elicitation price,
we cannot recover a complete demand curve at the individual level.
The resulting \texttt{wtp\_high} indicator conflates the household's
underlying valuation with the random price draw: a household assigned
a high price but that accepts may have different characteristics than one
assigned a low price and accepting. The median split on WTP-stated
is therefore a coarse proxy for true willingness-to-pay.

\paragraph{Null average effects constrain power.}
The original papers find small average effects on \texttt{life\_index}
($\text{ITT} \approx 0.09$ SD in our pooled sample). With such small
signal, detecting heterogeneity of plausible magnitude would require
a substantially larger sample or stronger variation in WTP.

\paragraph{D6 is exploratory and R2-only.}
The perceptions-of-security index (\texttt{r\_crime\_index}) is not
available in R1 and does not survive FDR correction in the original
paper ($q > 0.30$ in \citealt{lee2019jpe}). Results for D6 should be
treated as exploratory only.

\paragraph{Attrition.}
The \texttt{demand.dta}--\texttt{impacts\_2.dta} merge (2014--2017)
has not been validated against the original paper's attrition tables.
The WTP merge rate in the pooled sample is 99.5\,\%, suggesting minimal
non-response, but differential attrition by WTP group cannot be ruled
out without a formal balance check.

% ── References ───────────────────────────────────────────────────────────────
\clearpage
\begin{singlespace}
\begin{thebibliography}{9}

\bibitem[Benjamini and Hochberg(1995)]{benjamini1995}
Benjamini, Yoav, and Yosef Hochberg. 1995.
``Controlling the False Discovery Rate: A Practical and Powerful
Approach to Multiple Testing.''
\textit{Journal of the Royal Statistical Society, Series B}
57~(1): 289--300.

\bibitem[Lee, Miguel, and Wolfram(2019)]{lee2019jpe}
Lee, Kenneth, Edward Miguel, and Catherine Wolfram. 2019.
``Experimental Evidence on the Economics of Rural Electrification.''
\textit{Journal of Political Economy} 128~(4): 1523--1565.

\bibitem[Lee, Miguel, and Wolfram(2020)]{lee2019jep}
Lee, Kenneth, Edward Miguel, and Catherine Wolfram. 2020.
``Does Household Electrification Supercharge Economic Development?''
\textit{Journal of Economic Perspectives} 34~(1): 122--144.

\end{thebibliography}
\end{singlespace}

\end{document}
